\documentclass[10pt,a4paper]{amsart}
\usepackage[margin=19mm]{geometry}
\usepackage{amsmath,amssymb,mathtools}
\usepackage[scr=rsfs,cal=euler]{mathalfa}
\usepackage{xfrac}
\usepackage[unicode,hidelinks,bookmarks=false]{hyperref}
\newcommand{\ie}{i.e.\ }
\newcommand{\cf}{cf.\ }
\newcommand{\resp}{respectively\ }
\newcommand{\Subsection}{\subsection}
\providecommand{\nobreakdash}{\nobreak\hbox{-}}
\setlength{\emergencystretch}{2em}
\multlinegap0pt
\usepackage{bm}
\usepackage{bbm}
\usepackage{dsfont}
\usepackage{xspace}
\usepackage{cancel}
\usepackage{mathtools}
\usepackage{amsthm}
\makeatletter
\@ifundefined{theorem}{\newtheorem{theorem}{Theorem}}{}
\makeatother
\newcommand{\ev}[1]{\langle #1 \rangle}
\title[Quantized Orbital Angular Momentum from Discrete Chaotic Pha]{Quantized Orbital Angular Momentum from Discrete Chaotic Phase Surfaces}
\author{Netzer Moriya}
\date{Source: arXiv:2507.10229v1 (CC BY 4.0)}
\begin{document}
\maketitle

\noindent\emph{Source and licence.} Quantized Orbital Angular Momentum from Discrete Chaotic Phase Surfaces. Version arXiv:2507.10229v1 (CC BY 4.0), distributed under the Creative Commons Attribution 4.0 International licence, as recorded in the arXiv licence metadata for this exact version and shown on the abstract page https://arxiv.org/abs/2507.10229v1. Full rights notice: licenses/arxiv-2507.10229-CC-BY-4.0.txt.\par\smallskip

\noindent\textit{Editorial scope.} Counted unit: the Theorem whose source label is \texttt{theorem:OAM\_Selection\_Rule} in the article (source0.tex lines 108--111, source label \texttt{theorem:OAM\_Selection\_Rule}), delivered together with the article's own complete original proof of it (source0.tex lines 113--157, one \texttt{proof} environment of 45 lines). No mathematics is written, repaired or re-derived: statement and proof are the article's own text line for line. Required context retained uncounted: eq:ensemble\_average (lines 102--105). Every definition, display and label of those lines is retained unchanged. Omitted from this package and not redistributed: 1--101, 106--107, 112--112, 158--547. Nothing inside a retained excerpt is omitted or altered: every retained body is delivered exactly as the article's lines are written, so no omission receipt of any kind accompanies this package. Citations: the retained excerpts carry no citation command at all, so no bibliography entry of the article is transcribed and no reference is redistributed. Reference adaptation: none was needed and none was made; every reference inside the delivered unit resolves inside it, so no unresolved reference or citation is produced. Printed slips kept literal: no printed word, symbol, hypothesis or number of the counted statement or of its proof was altered. Layout and class: the article's own class is \texttt{article} with packages including amsmath, amsfonts, amsthm, amssymb, inputenc, fontenc, graphicx, caption, subcaption, hyperref, palatino, url, multicol, txfonts; the wrapper is the lane's pinned \texttt{amsart} editorial wrapper at 10pt on A4, which restates the theorem-like environments the retained text needs and adds the article's own macro definitions that the retained text needs (\texttt{\textbackslash ev}), with extra wrapper packages (bm, bbm, dsfont, xspace, cancel, mathtools). The delivered numbering is the unit's own. Assets: no retained span contains a figure environment, a graphics-inclusion command, a PDF-page inclusion, a table or a TikZ picture; no image file is copied into this package, the package declares no asset dependency and no image file is redistributed with it. Licence: version arXiv:2507.10229v1 of this article is distributed under the Creative Commons Attribution 4.0 International licence, as recorded in the arXiv licence metadata for this exact version and shown on the abstract page https://arxiv.org/abs/2507.10229v1. A full rights notice is delivered at licenses/arxiv-2507.10229-CC-BY-4.0.txt. Characters: every retained excerpt is pure ASCII, so the delivered engine, XeLaTeX with its default Latin Modern text font, reports no missing character.
\begin{equation}
    \ev{a_\ell(r)} = \frac{E_0}{2\pi} \int_0^{2\pi} \ev{e^{i\phi(r,\theta)}} e^{-i\ell\theta} d\theta.
    \label{eq:ensemble_average}
\end{equation}

\begin{theorem}[OAM Selection Rule]
Let $\phi(r,\theta)$ be a random phase surface and assume that $\ev{e^{i\phi(r,\theta)}}$ exists and is finite for all $(r,\theta)$. Then the ensemble-averaged OAM amplitude $\ev{a_\ell(r)}$ is non-zero if and only if the ensemble-averaged phase factor $\ev{e^{i\phi(r,\theta)}}$ possesses a non-zero Fourier coefficient at angular frequency $\ell \in \mathbb{Z}$.
\label{theorem:OAM_Selection_Rule}
\end{theorem}

\begin{proof}
From Eq.~\ref{eq:ensemble_average}, we have:
\begin{equation}
    \ev{a_\ell(r)} = \frac{E_0}{2\pi} \int_0^{2\pi} \ev{e^{i\phi(r,\theta)}} e^{-i\ell\theta} d\theta.
\end{equation}

Since $\ev{e^{i\phi(r,\theta)}}$ is $2\pi$-periodic in $\theta$ and bounded (as $|\ev{e^{i\phi(r,\theta)}}| \leq 1$), it admits a Fourier series representation:
\begin{equation}
    \ev{e^{i\phi(r,\theta)}} = \sum_{n=-\infty}^{\infty} C_n(r) e^{in\theta},
    \label{eq:fourier_expansion}
\end{equation}
where the Fourier coefficients are given by:
\begin{equation}
    C_n(r) = \frac{1}{2\pi} \int_0^{2\pi} \ev{e^{i\phi(r,\theta)}} e^{-in\theta} d\theta.
    \label{eq:fourier_coefficients}
\end{equation}

Substituting Eq.~\ref{eq:fourier_expansion} into the expression for $\ev{a_\ell(r)}$:
\begin{equation}
    \ev{a_\ell(r)} = \frac{E_0}{2\pi} \int_0^{2\pi} \left[\sum_{n=-\infty}^{\infty} C_n(r) e^{in\theta}\right] e^{-i\ell\theta} d\theta.
\end{equation}

Assuming uniform convergence of the Fourier series (which holds for bounded periodic functions), we may interchange the sum and integral:
\begin{equation}
    \ev{a_\ell(r)} = \frac{E_0}{2\pi} \sum_{n=-\infty}^{\infty} C_n(r) \int_0^{2\pi} e^{i(n-\ell)\theta} d\theta.
\end{equation}

Applying the orthogonality relation for complex exponentials:
\begin{equation}
    \int_0^{2\pi} e^{i(n-\ell)\theta} d\theta = 2\pi \delta_{n,\ell},
\end{equation}
where $\delta_{n,\ell}$ is the Kronecker delta, we obtain:
\begin{equation}
    \ev{a_\ell(r)} = E_0 C_\ell(r).
    \label{eq:final_result}
\end{equation}

Therefore:
\begin{itemize}
    \item $\ev{a_\ell(r)} \neq 0$ if and only if $C_\ell(r) \neq 0$
    \item By Eq.~\ref{eq:fourier_coefficients}, $C_\ell(r) \neq 0$ if and only if $\ev{e^{i\phi(r,\theta)}}$ has a non-zero Fourier component at frequency $\ell$
\end{itemize}

This completes the proof of the bidirectional statement.
\end{proof}

\end{document}
